\section{Discussion}\label{sec:discussion}

\subsection{What is new?}

The individual tools used by \cage are familiar: policy masks, gradient-boosted trees, quantile margins, and Jain fairness are established techniques. The contribution lies in how they are assigned distinct responsibilities in an online scheduler. Existing network-aware systems optimize communication or response time; policy frameworks enforce placement rules; temporal resource managers predict or react to SLO problems. \cage places these components in a strict order that reflects the semantics of sensitive workloads. First, legality is not a penalty. An illegal node is absent from the candidate set, so no performance score can compensate for a policy failure. Second, prediction is temporal and evaluated at the scheduler level. Current telemetry is a strong baseline, but the experiments show that history reduces migration and temporal violations under matched conditions. Third, fairness is implemented before final ranking rather than reported only after the fact. Fourth, the empirical projection is applied at commit time and is evaluated through the no-projection ablation. The result is a transparent pipeline in which every migration can be explained as policy-only, risk-only, projection-only, or a combination.

\subsection{Interpretation of the trade-off}

The complete method does not dominate every baseline on every metric. Against policy-Q95, it performs slightly more migrations to remove temporal violations and reduce the native unsafe proxy. Against a stricter 0.99 boundary and the current-HGB baseline, it performs fewer migrations while improving temporal compliance. Against the no-projection ablation, it performs about one percentage point more migrations to eliminate Q95 violations, with no change in the native proxy. These are useful and defensible trade-offs. Claiming universal superiority would hide the operational choice between migration cost and risk tolerance. The fairness results also require context. Least-loaded placement has very low Jain fairness because it repeatedly chooses the same apparent minimum. Random placement is more distributed but has worse temporal violations and higher proxy risk. The fair shortlist improves distribution without making randomness the objective. The balanced workload has lower absolute Jain values because it contains fewer events and destinations, so cross-workload fairness values should not be compared as if sample sizes were identical.

\subsection{Threats to validity}

\textbf{Construct validity.} The policy metadata are synthetic. They test whether the scheduler obeys hard rules and remains feasible across policy distributions; they do not describe the real regulatory distribution of Alibaba workloads. The native next-state unsafe metric is also a proxy. It indicates what happened naturally to the selected node, not what would have happened after receiving a migrated service.

\textbf{Internal validity.} Temporal leakage is controlled by constructing labels, thresholds, residual margins, and model parameters only from past or calibration data. Event detection uses the instance burst ending at $t$, node state is observed at $t$, and outcome is read at $t+30$ seconds. The final timestamp is excluded. Deterministic seeds and event ordering make the replay reproducible.

\textbf{External validity.} The evaluation uses one microservice-resource shard from one trace provider. Production clusters may have different telemetry scales, admission rules, actuation delays, and service graphs. The results therefore establish performance under the stated trace and replay protocol rather than universal generalization across platforms. The WSL2 timing study measures algorithmic overhead only; it does not represent end-to-end control-plane latency or migration completion time.

\textbf{Model validity.} Histogram gradient boosting performs well on this trace, but concept drift can change calibration. A deployment should monitor Brier score, calibration error, acceptance rate, and threshold-crossing misses, and should retrain only from temporally valid data. The empirical Q95 guardrail is not a proof of safety and does not account for candidate-specific interference, network staging, thermal constraints, or the transient old-plus-new resource demand of rolling migration.

\subsection{Practical deployment}

A Kubernetes implementation can map the four layers to familiar extension points. An admission controller or scheduler filter plugin enforces policy. A node-telemetry service maintains the three-snapshot feature state and risk scores. The shortlist and ranking execute in a scheduler score plugin or external scheduler. A reserve or pre-bind hook repeats the policy, temporal, and Q95 checks and records the destination reservation. If the recheck fails, the scheduler selects the next candidate or emits an explicit infeasibility decision. This structure avoids silently falling back to an illegal node. Policy changes and attestation revocation require immediate mask updates. A service whose current source becomes illegal should be migrated even if its operational risk is low. Conversely, a temporary load spike should not relax jurisdiction or TEE requirements. The separation of policy and risk makes these cases straightforward.
